\documentclass[11pt, a4paper]{article}

% --- Paquets ---
\usepackage[utf8]{inputenc}
\usepackage[T1]{fontenc}
\usepackage[french]{babel}
\usepackage[margin=2.45cm]{geometry}
\usepackage{hyperref}
\usepackage{enumitem}
\usepackage{booktabs}
\usepackage{tabularx}
\usepackage{graphicx}
\usepackage{float}
\usepackage{xcolor}

\hypersetup{
    colorlinks=true,
    linkcolor=blue,
    urlcolor=blue,
    pdftitle={Bilan hebdomadaire - Semaine du 15/06/2026}
}

\newcommand{\placeholderfigure}[2]{
\begin{figure}[H]
    \centering
    \fbox{\parbox{0.93\textwidth}{
        \vspace{1.2cm}
        \centering
        \textbf{Emplacement pour capture / graphique}\par
        \vspace{0.2cm}
        \texttt{#1}\par
        \vspace{0.2cm}
        #2
        \vspace{1.2cm}
    }}
    \caption{#2}
\end{figure}
}

\title{\textbf{Bilan hebdomadaire TER} \\ \large Semaine du 15 Juin 2026}
\author{ABBES Samy \& AREZKI Célia}
\date{}

\begin{document}

\maketitle

\section{Contexte général}
Cette semaine a été consacrée à la stabilisation finale de notre plateforme d'expérimentation ByzFL, à l'organisation des résultats produits par les différentes campagnes de tests, ainsi qu'à la mise en place d'une chaîne de comparaison en deux niveaux : d'abord comparer tous les agrégateurs pour chaque attaque, puis comparer uniquement les meilleures courbes retenues pour chaque attaque afin d'identifier l'attaque la plus difficile et le meilleur agrégateur associé.

Nous avons travaillé en binôme de manière continue sur la base de code, sur les scénarios JSON, sur les résultats générés et sur les scripts de comparaison. Cette semaine a aussi été marquée par plusieurs opérations Git importantes : pull, rebase, correction de conflit sur un CSV de résultats et push final vers le dépôt distant.

\section{Objectifs de la semaine}
Les objectifs fixés au début de cette semaine étaient les suivants :
\begin{itemize}
    \item ajouter de nouveaux scénarios et enrichir ceux déjà présents dans le dépôt ;
    \item organiser les scénarios par attaque et par distribution pour rendre l'arborescence plus lisible ;
    \item lancer les scénarios IID et Non-IID de façon séparée afin d'éviter les conflits de fichiers de résultats ;
    \item créer deux comparateurs complémentaires : un comparateur attaque/agrégateur et un comparateur final basé sur la meilleure courbe de chaque attaque ;
    \item sécuriser la synchronisation Git entre les résultats produits localement et les résultats du dépôt distant ;
    \item préparer une version exploitable dans le rapport hebdomadaire avec des emplacements pour les figures finales.
\end{itemize}

\section{Historique des modifications et commits}
Pour garder une trace claire du travail réalisé, nous avons repris l'historique Git de la semaine et identifié les changements les plus importants.

\subsection{Commit de structuration des scénarios non-IID}
Une première série de modifications a concerné la normalisation des configurations non-IID. Nous avons corrigé les scénarios pour homogénéiser les paramètres et garantir la cohérence des comparaisons, notamment en imposant une base commune de \textbf{5 clients honnêtes} et \textbf{2 attaquants} sur les scénarios concernés.

Cette étape a touché de nombreux fichiers JSON dans les dossiers des scénarios : SignFlipping, Inner Product Manipulation, Opt-IPM, ALIE, Opt-ALIE, Mimic, Infinity et Gaussian.

\subsection{Commit lié au lancement des expériences Mimic}
Un autre commit majeur de la semaine a ajouté et validé les scénarios Mimic en IID, avec tous les agrégateurs retenus. Cela a généré de nouveaux graphiques de sortie et a enrichi le fichier global de résultats `resultats_combines.csv`.

\subsection{Commit de modification du lanceur d'expériences}
Nous avons ensuite modifié le lanceur principal pour améliorer l'organisation des fichiers de sortie. Le but était d'éviter les collisions entre expériences IID et Non-IID, en attribuant des noms plus précis aux CSV générés selon l'attaque et la distribution.

\subsection{Commit des commandes Google Colab}
Enfin, un commit a porté sur le fichier de commandes Colab. Nous avons clarifié la procédure d'exécution GPU et la séquence Git à suivre lorsque les expériences sont relancées depuis Colab.

\section{Organisation du dépôt et des scénarios}
L'un des apports les plus importants de cette semaine a été l'organisation définitive des scénarios par attaque.

Chaque famille d'attaque est désormais rangée dans un sous-dossier dédié, avec deux sous-ensembles :
\begin{itemize}
    \item \texttt{iid/} pour les scénarios à distribution identique entre les clients ;
    \item \texttt{non\_iid/} pour les scénarios hétérogènes.
\end{itemize}

Cette nouvelle structure concerne notamment :
\begin{itemize}
    \item Sign Flipping,
    \item Inner Product Manipulation (IPM),
    \item Optimal Inner Product Manipulation (Opt-IPM),
    \item A Little Is Enough (ALIE),
    \item Optimal A Little Is Enough (Opt-ALIE),
    \item Mimic,
    \item Infinity,
    \item Gaussian.
\end{itemize}

Nous avons également harmonisé les noms des agrégateurs utilisés dans les scénarios : Average, Median, TrMean, Geometric Median, Krum, Multi-Krum, Centered Clipping, MDA, MoNNA, Meamed, CAF et SMEA.

\begin{figure}[H]
    \centering
    \includegraphics[width=0.9\linewidth]{CAP.png}
    \caption{Organisation finale des scénarios par attaque avec sous-dossiers IID et Non-IID}
    \label{fig:arborescence_scenarios}
\end{figure}

\section{Passage à deux niveaux de comparaison}
Cette semaine a surtout servi à mettre en place une logique d'analyse plus claire des résultats.

\subsection{Premier niveau : tous les agrégateurs pour chaque attaque}
Le premier script de comparaison, \texttt{comparer\_attaques\_agregateurs.py}, trace pour chaque attaque les courbes de tous les agrégateurs disponibles. Cela permet de visualiser la robustesse de chaque défense face à une attaque donnée, en gardant un axe des ordonnées normalisé entre 0 et 1 et des couleurs cohérentes d'un graphique à l'autre.

Le script produit également un résumé CSV des meilleurs agrégateurs par attaque, avec un classement interne selon l'accuracy finale.

\begin{figure}[H]
    \centering
    \includegraphics[width=0.9\linewidth]{MIMIC.png}
    \caption{Courbes de tous les agrégateurs pour une attaque donnée}
    \label{fig:comparaison_agregateurs}
\end{figure}

\subsection{Deuxième niveau : meilleure courbe retenue pour chaque attaque}
Le second script, \texttt{comparer\_meilleures\_attaques.py}, prend le meilleur agrégateur pour chaque attaque, extrait sa courbe complète et compare ensuite ces meilleures courbes sur un seul graphique final.

L'idée est de répondre à deux questions distinctes :
\begin{itemize}
    \item quelle est la pire attaque parmi celles testées ?
    \item quel est le meilleur agrégateur pour chaque attaque, puis globalement ?
\end{itemize}

Ce deuxième niveau fournit donc un classement final plus synthétique et plus exploitable pour le rapport.

\begin{figure}[H]
    \centering
    \includegraphics[width=0.9\linewidth]{courbe_comparaison_meilleures_attaques.png}
    \caption{Comparaison finale des meilleures courbes retenues pour chaque attaque}
    \label{fig:comparaison_meilleures_attaques}
\end{figure}

\section{Résultats observés cette semaine}
À la fin de cette semaine, les résultats consolidés ont montré que, sur les attaques effectivement présentes dans les CSV fusionnés, les meilleures performances finales atteignent des niveaux très élevés.

Dans la dernière comparaison finale générée, nous observons par exemple :
\begin{itemize}
    \item \textbf{SignFlipping} : meilleure courbe avec \textbf{Average}, accuracy finale d'environ \textbf{0.9772} ;
    \item \textbf{InnerProductManipulation} : meilleure courbe avec \textbf{MultiKrum}, accuracy finale d'environ \textbf{0.9811} ;
    \item \textbf{Infinity} : meilleure courbe avec \textbf{MultiKrum}, accuracy finale d'environ \textbf{0.9836}.
\end{itemize}

Ces résultats confirment que certains agrégateurs robustes restent très performants, même lorsque les attaques deviennent plus agressives. Ils montrent aussi que le comportement dépend fortement du couple \emph{attaque / agrégateur}.

\placeholderfigure{Tableau ou extrait CSV à insérer}{Extrait du classement final des attaques selon la meilleure accuracy finale}

\section{Résolution du conflit Git et gestion des résultats}
Cette semaine a aussi été importante du point de vue du versionnement.

Lors de nos synchronisations avec le dépôt distant, nous avons rencontré un cas concret de conflit sur le fichier généré \texttt{resultats\_experiences/resultats\_combines.csv}. Le problème ne venait pas du code source, mais du fait que ce fichier de résultats était modifié à la fois localement sur Colab et dans le dépôt distant.

Nous avons alors pris la décision de ne pas écraser les résultats, mais de conserver les deux jeux de résultats en les fusionnant. Cette approche nous a conduits à clarifier la procédure de travail sur Colab :
\begin{itemize}
    \item exécuter les expériences ;
    \item générer les CSV et les graphiques ;
    \item récupérer les changements distants avec \texttt{git pull --rebase} ;
    \item résoudre les conflits de résultats en fusionnant les lignes utiles du CSV ;
    \item relancer les scripts de comparaison si nécessaire ;
    \item terminer par un \texttt{git push} propre.
\end{itemize}

Nous avons aussi modifié le lanceur pour que les sorties puissent être distinguées plus facilement selon la famille exécutée, par exemple avec des noms de fichiers du type \texttt{resultats\_<attaque>\_<iid/non\_iid>.csv} lorsqu'on lance une série homogène de scénarios.

\placeholderfigure{Capture du conflit Git ou du terminal Colab}{Gestion d'un conflit de rebase sur le fichier de résultats et résolution manuelle}

\section{Travail d'intégration et de nettoyage}
En parallèle des expérimentations et de la gestion Git, nous avons poursuivi le travail de nettoyage du dépôt :
\begin{itemize}
    \item mise à jour des noms de scripts en français ;
    \item révision des chemins d'accès ;
    \item centralisation des comparaisons dans des dossiers dédiés ;
    \item mise en cohérence des CSV, images et sous-dossiers de sortie ;
    \item préparation du rapport hebdomadaire avec des emplacements dédiés pour les captures.
\end{itemize}

Ce travail d'organisation est important car il conditionne la lisibilité du projet final et la qualité de la démonstration devant les encadrants.

\section{Bilan technique}
Au terme de cette semaine, nous avons obtenu une base beaucoup plus solide qu'au début de la période :
\begin{itemize}
    \item les scénarios sont mieux rangés et mieux nommés ;
    \item les paramètres des expériences sont harmonisés ;
    \item les sorties de résultats sont plus sûres et plus lisibles ;
    \item la comparaison se fait maintenant en deux niveaux, ce qui clarifie l'analyse ;
    \item la synchronisation Git et Colab est mieux maîtrisée malgré les conflits sur les fichiers générés.
\end{itemize}

\section{Conclusion}
La semaine du 15/06 a été une semaine de consolidation et de structuration. Nous avons non seulement enrichi le contenu expérimental, mais aussi renforcé la qualité du pipeline complet : scénarios, exécution, comparaison, gestion des sorties et synchronisation avec le dépôt distant.

Le projet a maintenant une organisation plus professionnelle, avec des résultats mieux séparés selon la distribution et l'attaque, et avec des outils de comparaison capables d'identifier à la fois le meilleur agrégateur pour chaque attaque et l'attaque la plus difficile à défendre.

La prochaine étape consistera à lancer les nouvelles familles de scénarios en IID et Non-IID pour toutes les attaques, à remplir les résultats correspondants et à compléter la comparaison finale avec l'ensemble des attaques prévues.

\end{document}